% \subsection{Image Classification with CNNs}
% \label{sec:classification}

% We trained deep neural networks to perform classification on the CIFAR-10, CIFAR-100, and ImageNet datasets. We trained the networks with several optimization algorithms and their Lookahead counterparts.
% higher learning rates mean faster convergence but not necessarily better generalization
\subsection{CIFAR-10 and CIFAR-100} 
\label{sec:cifar}

The CIFAR-10 and CIFAR-100 datasets for classification consist of $32 \times 32$ color images, with 10 and 100 different classes, split into a training set with 50,000 images and a test set with 10,000 images. We ran all our CIFAR experiments with 3 seeds and trained for 200 epochs on a ResNet-18 \citep{he2016deep} with batches of 128 images and decay the learning rate by a factor of $5$ at the 60th, 120th, and 160th epochs.  Additional details are given in appendix \ref{app:experiments}.

We summarize our results in Figure \ref{tabel:cifar-optimizer-compare}.\footnote{We refer to SGD with heavy ball momentum \citep{polyak1964some} as SGD.} We also elaborate on how Lookahead contrasts with SWA and present results demonstrating lower validation error with Pre-ResNet-110 and Wide-ResNet-28-10 \cite{Zagoruyko2016WRN} on CIFAR-100 in appendix \ref{app:swa-comparison}.  Note that Lookahead achieves significantly faster convergence throughout training even though the learning rate schedule is optimized for the inner optimizer---future work can involve building a learning rate schedule for Lookahead. This improved convergence is important for better anytime performance in new datasets where hyperparameters and learning rate schedules are not well-calibrated.


\begin{figure}[t]
\centering
\begin{minipage}[t]{0.48 \linewidth}
    \centering
    \includegraphics[width=1. \textwidth]{images/CIFARTrainLoss.pdf}
\end{minipage}
\hfill
\begin{minipage}[t]{0.48 \linewidth}
\begin{table}[H]
\begin{center}
%\begin{small}
\vspace{-1.4in}
\begin{sc}
\begin{tabular}{l c c c c }
\toprule
Optimizer & CIFAR-10 & CIFAR-100  \\
\midrule
SGD & $95.23 \pm .19$ & $78.24 \pm .18$ \\ 
Polyak & $95.26 \pm .04$ & $77.99 \pm .42$ \\
Adam & $94.84 \pm .16$ & $76.88 \pm .39$ \\ 
Lookahead & $95.27 \pm .06$ & $78.34 \pm .05$ \\
\bottomrule
\end{tabular}
\end{sc}
%\end{small}
\end{center}
\caption{CIFAR Final Validation Accuracy.}
\end{table}
\end{minipage}
\caption{Performance comparison of the different optimization algorithms. ({\bf{Left}}) Train Loss on CIFAR-100. ({\bf{Right}}) CIFAR ResNet-18 validation accuracies with various optimizers. We do a grid search over learning rate and weight decay on the other optimizers (details in appendix \ref{app:experiments}). Lookahead and Polyak are wrapped around SGD.}
\label{tabel:cifar-optimizer-compare}
\end{figure}



% \begin{table}[t]
% \caption{CIFAR Test Accuracy for various $k$}
% \label{tabel:steps}
% \vskip 0.1in
% \begin{center}
% \begin{small}
% \begin{sc}
% \begin{tabular}{l c c}
% \toprule
% $k$ & CIFAR-10 & CIFAR-100 \\
% \midrule
% 2 & 95.14 & 76.99\\
% 3 & 95.20 & 76.47 \\ 
% 5 & 94.97 & 76.60\\ 
% 10 & 95.01 & 76.73 \\
% 20 &  95.21 &  77.57 \\
% \bottomrule
% \end{tabular}
% \end{sc}
% \end{small}
% \end{center}
% \vskip -0.1in
% \end{table}




